\documentclass[11pt,a4paper]{article}
\usepackage[utf8]{inputenc}
\usepackage[T1]{fontenc}       % Crucial for proper Czech/Slovak font rendering
\usepackage[czech]{babel}
\usepackage{amsmath, amssymb}
\usepackage{geometry}
\geometry{margin=2.5cm}

\begin{document}

\title{Seznam požadavků ke zkoušce z Datových struktur II\\-- Kapitoly a pojmy k ověření --}
\author{Matfyz}
\date{\today}
\maketitle

\section{Statické slovníky}

\subsection{Teoretický úvod a výpočetní model}
\begin{itemize}
    \item \textbf{Výpočetní model Word-RAM}:
    \begin{itemize}
        \item Definice velikosti slova $w$ a vztah k velikosti univerza $U$ ($u = 2^w$).
        \item Přehled operací prováděných v konstantním čase $O(1)$.
        \item Paměťová omezení a adresace buněk.
    \end{itemize}
    \item \textbf{Formulace statického slovníkového problému}:
    \begin{itemize}
        \item Časový invariant operace \texttt{Lookup} v nejhorším případě ($O(1)$).
        \item Prostorový požadavek na striktně lineární paměť ($O(n)$).
    \end{itemize}
\end{itemize}

\subsection{Perfektní hešování FKS (Fredman, Komlós, Szemerédi)}
\begin{itemize}
    \item \textbf{Definice $c$-univerzálního systému}:
    \item \textbf{Lemma o střední hodnotě počtu kolizí, při výběru náhodné heš. funkce}:
    \begin{itemize}
        \item Očekávaný počet kolizí v závislosti na parametrech $c$, $n$ a $m$.
    \end{itemize}
    \item \textbf{Nalezení perfektní hešovací funkce h pro $ m = \left\lceil n^2 \cdot c \right\rceil$ a $ m = n $ }:
    \begin{itemize}
        \item Markovova nerovnost.
        \item Lemma o džbánu.
        \item $\mathbb{E}[\# \textit{pokusů než najdeme takovou h}]$.
    \end{itemize}
    \item \textbf{FKS algoritmus}:
    \begin{itemize}
        \item Dvouúrovňové hešování: $ f, g_i $.
        \item \textit{Lookup(x)}.
        \item Velikost všech $Q_i$ je $ O(n) $ (trik $ p_i^2 = p_i + 2\binom{p_i}{2} $).
        \item Spotřeba paměti a čas na konstrukci.
    \end{itemize}
    \item \textbf{Odbočka o třídění $ x \in_R [0,1] $}:
    \begin{itemize}
        \item $ n $ přihrádek.
        \item Bubblesort v přihrádce v $ O(p_i^2) $.
        \item Proč to má $ O(n) $ složitost?
    \end{itemize}
\end{itemize}
\subsection{Deterministické slovníky}
\begin{itemize}
    \item \textbf{Redukce z obecného univerza na $ O(n^2) $.}
    \item \textbf{Definice $q$-dobré funkce}
    \item \textbf{Lemma pokud $(f,g)$ je $q$-dobrá a $r \geq \log(n) + 1$, pak $ \exists (f', g') $, která je $q'$-dobrá}
    \begin{itemize}
        \item Umím vyrobit z funkce $(f,g)$ novou, s méně kolizema.
    \end{itemize}
    \item \textbf{Tříkrokový algoritmus, pro nalezení perfektní hešovací funkce}:
    \begin{itemize}
        \item Zvolit $r$. Množina $S \subset {0,1}^w$, kde $w \leq 2\log(n) + O(1)$.
        \item $(f,g) \rightarrow (f',g') \rightarrow (f'',g'')$.
        \item $(f'', g'')$ je \textit{0-dobrá} funkce.
    \end{itemize}
    \item \textbf{Zahešování $x$}:
    \begin{itemize}
        \item Čas $O(1)$ a prostor $O(n)$.
    \end{itemize}
    \item \textbf{Důkaz Lemmatu}:
    \begin{itemize}
        \item Seřazení řádků.
        \item Volíme posuny $a_i$ náhodně, odhadneme $\#$ nových kolizí ($\#NK$) nově posunutého řádku $S_i$ s do teď posunutými řádky $|S_{<i}|$.
        \item Zkoušíme trefit $a_i$ do $O(1)$ pokusů (třeba 2). Ukážeme, kolik maximálně bude střední hodnota $\#NK$. 
        \item Všechny $a_i$ jdou najít efektivně v čase průměrně $O(n)$.
        \item Odhadneme počet všech kolizí, přes všechny řádky.
        \begin{itemize}
            \item Vyčistíme sumu od konce: triviální případy $|S_i| = 0$ nebo $|S_i| =1$.
            \item Omezíme křížové součiny $|S_i|\cdot|S_{<i}|$.
            \item Využijeme odhadu $d^2 \leq 4 \cdot \binom{d}{2}$, což platí pouze pro nějaká $d$, k provázání $|S_j|^2$ s $q$, kdy $q$ je celkový počet kolizí původních.
            \item Dosadíme do velké původní sumy pro $\#$\textit{kolizí} a dostaneme odhad $\#$\textit{kolizí} $\leq n \cdot \left\lfloor 2^{3-r} \cdot q \right\rfloor $. 
        \end{itemize}
    \end{itemize}
    \item \textbf{Derandomizace}:
    \begin{itemize}
        \item Metoda podmíněné střední hodnoty.
        \item Jak počítat bity $a_i = {0,1}^r$ efektivně?
        \item Složitost zpracování jednoho řádku, plus složitost datové struktury, ve které \\ si udržujeme počty prvků, v sloupečcích a podle ní počítáme bity $a_i$, pro řádek $S_i$.
        \item Celková složitost derandomizovaného algoritmu $O(n \cdot \log(n))$.
    \end{itemize}
\end{itemize}
\section{Celočíselné datové struktury}
\subsection{van Emde-Boasovy stromy}
\begin{itemize}
    \item \textbf{Anatomie vEB stromu}:
    \begin{itemize}
        \item $min(S)$
        \item $max(S')$
        \item Přihrádky $P_0,\ldots,P_{\sqrt{U}-1}$
        \item Sumární strom ... $Sum$
    \end{itemize}
    \item \textbf{Find($x$)}
    \item \textbf{Succ($x$) (nejbližší následník)}:
    \begin{itemize}
        \item Rychlá kontrola s $min$.
        \item $x = (i,j)$.
        \item Zkusíme přihrádku $P_i$.
        \item Jinak vrátíme $P_{Sum.Succ(i)}.min$ . 
    \end{itemize}
    \item \textbf{Insert($x$)}:
    \begin{itemize}
        \item Rychlá kontrola s $min$.
        \item Pokud neskončíme, vezmeme $x$.
        \item Upravíme $max$ podle $x$.
        \item $x = (i,j)$.
        \item Rekurzíme se podle:
        \begin{itemize}
            \item Pokud $P_i$ existuje\dots
            \item Pokud $P_i$ prázdná\dots
        \end{itemize}
    \end{itemize}
    \item \textbf{Delete($x$)}:
    \begin{itemize}
        \item Pokud prázdný strom.
        \item Jinak pokud mažeme minimum, tak se může změnit $x$\dots.
        \item Pak $x = (i,j)$.
        \item Zavoláme $P_i.Delete(j)$.
        \item Pokud se $P_i$ vyprázdnila, musíme aktualizovat $Sum$.
        \item Aktualizovat $max$, kdybychom ho smazali.
    \end{itemize}
    \item \textbf{Problém s velikostí potřebné paměti $O(|U|)$, kde $|U| = 2^w$}:
    \begin{itemize}
        \item Naivní řešení s obyčejným polem.
        \item Trik pomocí $M$, $I$, $X$.
        \item Operace $Read(i)$ a $Write(i, y)$.
    \end{itemize}


\end{itemize}

\subsection{Van-Emde Boasov "intervalový" strom}
    \begin{itemize}
        \item \textbf{Struktura stromu}:
        \begin{itemize}
            \item Charakteristický vektor $S$.
            \item $O(|U|)$ vrcholů.
            \item Hloubka $O(\log|U|) = O(w)$.
            \item Monotónnost cesty z listu do kořene.
        \end{itemize}
        \item \textbf{Operace}:
        \begin{itemize}
            \item Rychlost operací $Pred(x)$, $Succ(x)$, $Find(x)$, pokud budeme umět rychle bin. vyhledávat na cestě z listu do kořene, bude $O(LL(|U|))$, jinak $O(L|U|)$.
            \item Plus musíme udržovat $min$ a $max$ podstromů a spojový seznam prvků $S$, abychom zvládli i situaci, kdy se ptáme na $Pred$/$Succ$ prvku $x$, který je v množině $S$.
            \item Problém s $Updaty$, běží v $O(L|U|)$.
        \end{itemize}
    \end{itemize}
\subsection{X-fast trie}
    \begin{itemize}
        \item \textbf{Struktura stromu}:
        \begin{itemize}
            \item Vychází z "intervalového" stromu + kukaččí hešování.
            \item Ukládáme pouze vrcholy, které jsou potřeba, tedy vrcholy odpovídající prefixům prvků, které jsou v $S$.
            \item Vyhledáváme na cestě mezi listem a kořenem binárně, pomocí dotazu v $O(1)$, díky hešování.
            \item Místo neexistujících synů nějakého vrcholu, si uložíme odkaz \\ na levého/pravého následníka (list) v jeho podstromě, který existuje. 
        \end{itemize}
        \item \textbf{Operace}:
        \begin{itemize}
            \item Operace $Find$, $Pred$, $Succ$ běží v čase $O(LL(|U|))$.
            \item Operace $Insert$, $Delete$, běží v čase $O(L(|U|))$ průměrně amortizovaně.
        \end{itemize}
        \item \textbf{Paměťová složitost $O(n \cdot \log(|U|))$}
    \end{itemize}
\subsection{Y-fast trie}
    \begin{itemize}
        \item \textbf{Struktura stromu}
        \begin{itemize}
            \item Setříděné prvky rozdělíme na bloky velikosti $B = \Theta(\log (|U|))$.
            \item Každý blok svého reprezentanta (typicky největší $x$). Uvnitř bloků BVS \\ $\rightarrow$ operace v $O(L(|B|)) = O(LL(|U|))$.
            \item Globální úroveň je x-fast trie a do ní ukládáme pouze reprezentanty těchto BVS.
            \item Prvků v x-fast trii je $O(n/\log(|U|))$.
        \end{itemize}
        \item \textbf{Paměťová složitost}:
        \begin{itemize}
            \item Všechny BVS dohromady $O(n)$.
            \item Globální x-fast trie $O(\#\text{prvků} \cdot \log(|U|)) = O(n)$.
        \end{itemize}
        \item \textbf{Operace}:
        \begin{itemize}
            \item $Find(x)$, $Succ(x)$, $Pred(x)$ - globální strom mi řekne následníka a předchůdce, tedy dva BVS a v nich opakuji dotaz na $x$.
            \item Složitost $O(L(B)) = O(LL(|U|))$.
        \end{itemize}
        \item \textbf{Updaty}:
        \begin{itemize}
            \item Normálně probíhají, bez reorganizace globální x-fast trie
            \item Udržujeme invariant hustoty bloku $[\frac{1}{4}B,1B]$. Jakmile blok podteče/přeteče tak \\ slučujeme/dělíme.
            \item Amortizovaně složitost $O(1)$. Dokazujeme skrz složitost globální x-fast trie vs. jak často se to děje?
            \item \textbf{POZOR} na jednu věc u mazání. 
        \end{itemize}
    \end{itemize}
\section{Haldy}
\subsection{Srovnání hald}
\begin{itemize}
    \item \textbf{Binární vs. k-regulární halda}
    \begin{itemize}
        \item Složitosti operací $Insert$, $Min$, $ExtractMin$, $Decrease$, $Delete$, $Build$, $Merge$.
        \item Dijkstruv algo. používá $n \times Insert$ a $n \times ExtractMin$ a $m \times Decrease$.
        \item Pro $k$-regulární haldu, celková složitost minimální pro $k := \frac{m}{n}$.
        \item Binomiální umí navíc $Merge$ v $O(\log(n))$.
    \end{itemize}
\end{itemize}
\subsection{Binomiální halda}
\begin{itemize}
    \item \textbf{Binomiální strom $B_k$}:
    \begin{itemize}
        \item Dvě rekurzivní definice.
        \item Pozorování na počet vrcholů, počet hladin a maximální stupeň kořenu, což je i maximální stupeň ve stromě. 
    \end{itemize}
    \item \textbf{Binomiální halda}:
    \begin{itemize}
        \item Definice: Les binomiálních stromů, v každém vrcholu jeden prvek, platí haldové uspořádání.
        \item Počtem prvků je jednoznačně určeno, které binomiální stromy v posloupnosti existují.
        \item Do $O(\log (n))$ se vleze:
        \begin{itemize}
            \item řády stromů,
            \item $\#$ stromů v lese,
            \item hloubky,
            \item maximální $\#$ dětí.
        \end{itemize}
        \item Uložení v paměti, každý vrchol:
        \begin{itemize}
            \item prvek + priorita,
            \item řád (počet dětí),
            \item ukazatel na rodiče,
            \item ukazatel na nejlevější dítě,
            \item ukazatel na levého/pravého sourozence (zapojeni cyklicky, pokud jsou to kořeny tak\dots)
        \end{itemize}
    \end{itemize}
    \item \textbf{Operace}:
    \begin{itemize}
        \item $Min$ v $O(\log(n))$
        \item $Merge$ v $O(\log(n))$
        \item $Insert$ v $O(\log(n))$
        \item $Build$ v $O(n)$
        \item $ExtractMin$ v $O(\log(n))$
        \item $DecreaseKey$ v $O(\log(n))$
        \item $IncreaseKey$ v $O(\log(n))$
        \item $Delete$ v $O(\log(n))$
    \end{itemize}
\end{itemize}
\subsection{Líná binomiální halda}
    \begin{itemize}
        \item \textbf{Myšlenka lenosti.}
        \item \textbf{Složitost operací $Merge$, $Insert$ a $Build$.}
        \item \textbf{Analýza $ExtractMin$, složitost $O(t+\log(n))$, kde $t := \# \text{stromů v lese}$.}
        \begin{itemize}
            \item Správně si zvolit potenciál.
        \end{itemize}
    \end{itemize}
\subsection{Fibonacciho Halda}
    \begin{itemize}
        \item \textbf{Struktura haldy "Jak funguje, černý a bílý vrchol."}
        \item \textbf{Operace $Insert$, $Merge$, $Build$.}
        \item \textbf{Lemma 0 o Fibonacciho číslech}:
        \begin{itemize}
            \item Rekurzivní definice Fibonacciho čísla.
            \item Jiná přímá definice Fibonacciho čísla, kterou uvádí Lemma.
            \item Fibonacciho číslo roste exponenciálně!
        \end{itemize}
        \item \textbf{Lemma* o Tloušťce stromu}:
        \begin{itemize}
            \item Vrchol s $k$ dětmi má\dots
            \item Pomocné tvrzení v důkazu o řádech dětí nějakého vrcholu: $\forall i: r_i \geq i - 2$.
            \begin{itemize}
                \item \dots
            \end{itemize}
            \item Důkaz Lemma*, které
        \end{itemize}
        \item \textbf{Operace $Decrease$ a $Cut$.}
        \item \textbf{Těžká práce u $ExtractMin$.}
    \end{itemize}
\section{Splay stromy}
\subsection{Analýza Splay stromů}
    \begin{itemize}
        \item \textbf{Opakování Splay Stromů}
        \item \textbf{Věta Access Time}:
        \begin{itemize}
            \item Zadefinování pojmů jako $T(v)$, $s(v)$, $r(v)$, $\Phi(v)$.
            \item Znění: Amortizovaná cena $Splay(x)\leq 3(r'(x)-r(x))+1$, což je v $O(\log(n))$.
            \item Přepis $A_i = R_i + (\Phi_i -\Phi_{i-1})$ pro $i\in[m]$, za použití teleskopické sumy, kdy chceme odhadnout $R_i$. 
        \end{itemize}
        \item \textbf{Věta $m \times Splay$ trvá $O(m\cdot \log(n) + n \cdot \log(n))$ \dots} 
        \item \textbf{Váhová analýza}:
        \begin{itemize}
            \item Definice vah, pro vrcholy $w(v)$ a nové velikosti $s(v)$.
            \item Pozorování 1, že věta Access Time stále platí.
            \item Pozorování 2, $3(r'(x) -r(x))+1$ se už nevejde do $ O(\log(n))$, ale je potřeba lépe odhadnout, pomocí: $r_{max}$ a $r_{min}$.
        \end{itemize}
        \item \textbf{Analýza operací $Find(x)$, $Insert(x)$ a $Delete(x)$}
        \begin{itemize}
            \item Všechny se vlezou do $O(r_{max}-r_{min} + 1)$.
        \end{itemize}
        \item \textbf{Lemma $W$}"
        \begin{itemize}
            \item Vzorec pro to, jak dlouho trvá $Splay(x)$, pokud známe váhu $x$ a celkovou váhu stromu.
            \item Důkaz odhadneme $O(r_{max}-r_{min} + 1)$.
        \end{itemize}
        \item \textbf{Váhovou teorii s Lemmatem W aplikujeme na úplně obyčejný Splay strom, kde žádný prvek neupřednostňujeme, vyjde nám ta stará známá složitost $\mathcal{O}(\log n)$?}
        \begin{itemize}
            \item Nastavení vah rovnoměrně náhodně.
            \item Dosazení do Lemmat.
            \item Celková cena $m$ operací.
            \item Proč můžeme předpokládat, že $W=1$?
        \end{itemize}
        \item \textbf{Staticky optimální strom.}
        \item \textbf{Věta o Statické optimalitě}
        \begin{itemize}
            \item Znění, že Splay strom zpracuje asymptoticky sekvenci dotazů..
            \item Důkaz
            \begin{itemize}
                \item $w(x)=3^{-c_T(x)}$
                \item Celková váha $W<1$
                \item Dosazení do Lemmatu $W$
                \item Vyúčtování sekvence $m$ operací, tedy $\Phi_0 - \Phi_m$.
            \end{itemize}
        \end{itemize}
        \item \textbf{Static Finger bound \dots PRESKAKUJI}
        \item \textbf{Working Set Bound}
        \begin{itemize}
            \item Co to je Working set a velikost této množiny $z_i$.
            \item Znění věty.
            \item LRU seznam a že pozice v něm prvku $x$ je $p(x) = z(x)$.
            \item Nastavíme váhy $w(x) = \frac{1}{(p(x)+1)^2}$.
            \item Ukážeme, že $W$, je opět omezeno konstantou, tedy v $O(1)$.
            \item Vyúčtování $\Phi_0 - \Phi_m$ - NECHAPU.
        \end{itemize}
    \end{itemize}
\section{Cestové dotazy a updaty na stromech}
    \subsection{Cesta}
    \begin{itemize}
        \item \textbf{Motivace - kostry!}
        \item \textbf{Intervalový strom nad cestou}
        \item \textbf{Všechny cestové dotazy i updaty v $O(\log(n))$, pomocí kanonických podstromů a líného propagování změny.}
    \end{itemize}
    \subsection{Heavy-Light dekompozice (Strom)}
    \begin{itemize}
        \item \textbf{Definice HLD, definice těžké hrany.}
        \item \textbf{3 Pozorování}:
        \begin{itemize}
            \item Každý vrchol má maximálně 1 těžkou hranu
            \item Z toho plyne, že těžké hrany tvoří disjunktní cesty $\rightarrow$ můžeme na ně použít cestovou dekompozici.
            \item Každá cesta kořen-list obsahuje maximálně $O(\log(n))$ hran.
        \end{itemize}
        \item \textbf{Konstrukce pomocí dvou DFS v $O(n)$.}
        \item \textbf{LCA v HLD v $O(\log(n))$}
        \begin{itemize}
            \item Projdu po lehkých hranách.
            \item Pokud narazím na stejnou těžkou cestu $\rightarrow$, tak beru z těch dvou vyšších vrcholů, což v $O(1)$, díky tomu, co si každá těžká cesta pamatuje.
        \end{itemize}
        \item \textbf{Cestový dotaz/update}
        \begin{itemize}
            \item Najdeme $LCA(u,v)$, a rozdělíme dotaz na dvě cesty.
            \item Na každé navštívime max $O(\log(n))$ lehkých hran a těžkých cest
            \item Každý dotaz na těžké cestě v $O(\log(n))$, celkově tedy $O(\log^2(n) + \log(n)) = O(\log^2(n))$.
        \end{itemize}
    \end{itemize}
    \subsection{Link-Cut Stromy}
    \begin{itemize}
    \item \textbf{Struktura LCS a rozhraní}:
    \begin{itemize}
        \item Tlusté a tenké hrany.
        \item Strukturální operace: $Root(v)$, $Parent(v)$, $Cut(v)$, $Link(u, v)$, $Evert(v)$.
        \item Hodnotové operace: $Cost(v)$, $PathMin(v)$, $SetCost(v)$, $PathUpdate(v, \delta)$.
    \end{itemize}
    \item \textbf{Mikro-architektura: Splay Strom pro jednu tlustou cestu}
        \begin{itemize}
            \item Splay strom = tlustá cesta.
            \item Inorder průchod SS je pořadí na TC ze spoda nahoru.
            \item Líné vyhodnocování: vrcholy SS si udržují minimum podstromu, $+\delta$ a $ReverseFlag$
        \end{itemize}
    \item \textbf{Makro-architektura}
    \begin{itemize}
        \item Kořen Splay stromu $\neq$ nejvyšší vrchol na příslušné tlusté cestě v LCS.
        \item Kořen SS si pamatuje odkaz na svého lehkého rodiče.
        \item Jeden LCS je jeden ze stromů, který se snažíme reprezentovat.
    \end{itemize}
    \item \textbf{Líné vyhodnocování}
    \begin{itemize}
        \item Při průchodu zespoda nahoru vytvoříme zásobník, pak už jen projdeme zásobník a tím procházíme od shora dolů.
    \end{itemize}
    \item \textbf{Globální pohled + Operace $Expose(v)$}
    \item \textbf{Jak funguje $Expose(v)$ + amortizace\dots strany 39-41}
    \item \textbf{Aplikace na Dinicově algoritmu. - PRESKAKUJI}
    \end{itemize}
\section{Dynamizace}
\subsection{Dynamizace dvojrozměrných intervalových stromů}
\begin{itemize}
    \item \textbf{Dvojrozměrný intervalový strom}
    \begin{itemize}
        \item Struktura stromu, primární strom a sekundární stromy.
        \item Obdelníkové dotazy v čase $O(\log^2(n))$.
        \item Velikost stromu je pouze $O(n\cdot \log(n))$.
        \item Budování stromu v čase $O(n\cdot \log(n))$.
    \end{itemize}
    \item \textbf{Líné vyvážování}
    \begin{itemize}
        \item Vložení bodu.
        \item Udržujeme $\alpha$ vyváženost.
        \item Částečný rebuild, jakmile je někde podmínka porušena.
        \item Amortizovaná cena $\log^2(n)$.
    \end{itemize}
\end{itemize}
\subsection{Obecná dynamizace}
\begin{itemize}
    \item \textbf{Definice: vyhledávací problém}
    \begin{itemize}
        \item $f : U_Q \times P_{fin}(U_X) \rightarrow U_R$
    \end{itemize}
    \item \textbf{Rozložitelnost: klíčová vlastnost}
    \begin{itemize}
        \item $\exists \sqcup: U_R \times U_R \rightarrow U_R$, vyčíslitelná v konstatním čase, taková, že\dots
        \item Tři příklady rozložitelných a nerozložitelných problémů.
    \end{itemize}
    \item \textbf{Výchozí bod - Statická struktura, až SemiDynamizace a Dynamizace}
    \begin{itemize}
        \item Statická struktura - $B_S(n), Q_S(n), S_S(n)$ + rozumné chování těchto složitostí.
        \item Semidynamická struktura - $Q_D(n)$ (w.c.),$S_D(n)$, $\overline{I_D(n)}$ (amortizovaně). Pokud chceme dynamickou, ještě musíme přidat delete v čase $\overline{D_D(n)}$ amortizovaně.
        \item Struktura zdynamizování Statické.
        \item $Query$ - zeptáme se každého bloku a zkombinujeme a jeho $Q_D(n)$.
        \item $Insert$ - Najdeme nejbližší prázdný blok $B_i$, rozložíme bloky $<i$. Složíme do bloku $B_i$.
        \item Odvození, že $S_D(n) \leq S_S(n)$.
        \item Důkaz amortizované složitosti $\overline{I_D(n)}$.
    \end{itemize}
    \item \textbf{Worst-Case Dynamizace (STRANA 49)}
    \begin{itemize}
        \item Zlepšení těch dlouhých operací $Insert$, kdy dojde k přestavbě.
        \item Životní cyklus bloků, vždy žijou dva bloky od jednoho řádu $i$ a building bloku $i+1$ delame prubezne.
        \item Dotaz stále v $Q_D(n) \in O(Q_S(n) \cdot \log n)$
        \item Vkládání má novou w.c. složitost $I_D(n) \in O(\frac{B_S(n)}{n} \cdot \log n)$.
    \end{itemize}
    \item \textbf{Úplná dynamizace $\rightarrow$ přidání $Delete$}
    \begin{itemize}
        \item Pomocný adresář, pro pamatování, ve kterém bloku prvek je.
        \item Místo mazání prvku děláme pomníčky.
        \item Jakmile moc pomníčků $\rightarrow$ re-build.
        \item Amortizovaná cena $\overline{D_D(n)}$ je $O(\frac{B_S(n)}{n} \cdot \log n)$
    \end{itemize}
\end{itemize}

\section{Perzistence}
\subsection{Typy Perzistence}
\begin{itemize}
    \item \textbf{Efemérní, Semiperzistentní, Plně perzistentní, Funkcionální stuktury.}
    \item \textbf{Funkcionální seznamy s prependem - STRANA 51}
    \begin{itemize}
        \item Jednosměrný spojový seznam s přidáváním na začátek.
    \end{itemize}
    \item \textbf{Funkcionální BVS - STRANA 51}
    \begin{itemize}
        \item Kopírujeme cestu do kořene, verze je prostě odkaz na nový kořen.
        \item Paměť na verzi $O(\log n)$ a čas $O(\log n)$.
    \end{itemize}
    \item \textbf{Aplikace funkcionálních BVS - Lokalizace bodu v rovině rozdělené polygony STRANA 52 - 53}
    \begin{itemize}
        \item První řešení, pomocí dvou seřazených struktur, podle pásů a úseček v pásu $\rightarrow$ prostor $\in O(n^2)$, vybudování $\in O(n^2)$, dotaz $\in O(\log n)$.
        \item Lepší řešení pomocí funkcionálního BVS, kdy zametáme rovinu zeshora a vždy událost je přidání, odebrání úsečky. Pro každou událost postavíme verzi a každý pás úseček je pak tento funkc. BVS $\rightarrow$ dotaz stále $\in O(\log n)$, prostor $\in O(n \log n)$, vybudování $\in O(n \log n)$.
    \end{itemize}
\end{itemize}
\subsection{Semi-Persistentní DS}
\begin{itemize}
    \item \textbf{Pointerové DS}
    \begin{itemize}
        \item Krabička - má $O(1)$ dat, $O(1)$ ukazatelů. Držátko $O(1)$ ukazatelů.
        \item Operace co umíme s takovou DS: čtení krabičky, skok na jinou krabičku, modifikace krabičky, zakládání krabičky.
    \end{itemize}
    \item \textbf{Co vše musíme verzovat? Strukturální změny vs nestrukturální}
    \item \textbf{První technika: Tlusté vrcholy}
    \begin{itemize}
        \item Vrchol si pamatuje celou svoji historii.
        \item (počáteční vrchol, verze, změny)
        \item Update v $O(1)$. Ale omezený počet verzí.
        \item Query se zpomalí na $O(\log^2 n)$, kvůli binárnímu vyhledávání správné verze vrcholu.
    \end{itemize}
    \item \textbf{Obecná semi-perzistentní DS}
    \begin{itemize}
        \item Graf má omezený vstupní vrchol $\leq p$.
        \item Anatomie krabičky: Počáteční verze, data, žurnál ($e$ slotů změn, platí $p\leq e$), zpětné pointery.
        \item $Query$ - přijdu do krabičky, aplikuji všechny změny až do aktuální verze, těch je max $e$, což znamená, že konstatntně mnoho, tedy v $O(1)$.
        \item $Update$ - pokud volno, prostě připíšeme změnu a jinak musíme přetečenou krabičku kopírovat a musíme změnit zpětné pointery rodičů na tuhle aktuální krabičku.
        \item Trik na zacyklení, abychom nekopírovali do nekonečna.
        \item Důkaz amortizace, že každá strukturální změna se děje v $O(1)$ amortizovaně.
    \end{itemize}
    \item \textbf{Úplná perzistence DS}
    \begin{itemize}
        \item Problém s amortizovanou složitostí, můžeme se totiž vracet do verze, kdy došlo k drahému kopírování a vždy ho vynutit.
        \item Problém s porovnáním verzí, což se dá řešit projitím stromu "kolem", tedy linearizace historie.
        \item Musíme navíc omezit konstantou výstupní stupeň.
        \item Zbytek složitý\dots
    \end{itemize}
\end{itemize}
\section{Cache struktury}
\subsection{List Labeling problem}
\begin{itemize}
    \item \textbf{Definice problému. Jinými slovy List Order.}
    \item \textbf{LL v exponenciálním rozsahu}
    \begin{itemize}
        \item \# prvků $\leq M$.
        \item Zarážky jsou na pozicích $0, 2^M$.
        \item Nový prvek dám na průměr sousedů.
        \item Po $M$ insertech mezery aspoň 1.
        \item $O(1)$ na insert w.c.
        \item Funguje pouze pro $M \in O(w)$, kde $w$ je šířka slova.
    \end{itemize}
    \item \textbf{LL v polynomiálním rozsahu}
    \begin{itemize}
        \item Používáme bin. strom hloubky $O(\log M)$.
        \item Externí vrcholy jsou prvky.
        \item Kroky ve stromě L/P jsou bity v labelu.
        \item Dřivější prvky dostanou nižší labely.
        \item \# bitů labelu je $O(\log M)$.
        \item $Insert$ znamená změnit příslušný externí vrchol na interní a pod něj připojit přelabelovaný prvek a nový prvek.
        \item Vyvažovat budeme líně s BB-$\alpha$ stromy.
        \item Přebudování stojí $O(k)$, kde $k$ je \# prvků podstromu a musíme je přelabelovat, to se ale schová do $O(k)$ a tedy amortizovaně insert běží jako v líně vyvažovaném, tedy $O(\log n)$. 
    \end{itemize}
    \item \textbf{Finální List Order v $O(1)$ amortizovaně}
    \begin{itemize}
        \item Rozdělíme posloupnost na bloky velikosti $\theta(\log n)$
        \item Uvnitř bloků používáme labelování v exponenciálním rozashu.
        \item Nad bloky používáme labelování v poly.
        \item Prvek má (label bloku, label uvnitř bloku).
        \item Dokud $Insert$ do bloku, tak v $O(1)$.
        \item Jakmile přeteče, tak ho rozdělíme na dva poloviční a přidáme blok v globálních blocích, tedy musíme přelabelovat všechno uvnitř jednoho bloku.
        \item Čas $O(\log n)$ na přelabelovaní, čas na globální $Insert$ $O(\log n)$ amortizovaně. Amortizovaně $O(\frac{1}{\log n})$ krát.
        \item ZRADA, musí se přidat indirekce, kvůli přelabelování ostatních bloků při velkém globálním rebuildu, ve kterých je $O(\log n)$ prvků.  
    \end{itemize}
\end{itemize}

\end{document}